% =============================================================================
% NO2 en España — presentación (20 minutos)
% Compilar:  pdflatex slides.tex ; pdflatex slides.tex
%
% GUION TEMPORAL — 42 páginas en el cuerpo + 7 de reserva = 49
%
% ORDEN DE LAS TRES RAMAS: geoestadística -> areal -> procesos puntuales,
% que es el orden en que el curso las enumera. El acoplamiento va después,
% como síntesis. El resultado nulo areal deja de ser un desenlace y pasa a ser
% la PREGUNTA que responden las partes 5 y 6.
% -----------------------------------------------------------------------------
%   Portada e índice ...................... 0:40   (2 páginas)
%   1. El problema ........................ 2:20   (1 divisoria + 3)
%   2. Datos y diseño ..................... 2:10   (1 + 3)
%   3. El campo de fondo .................. 3:40   (1 + 5)
%   4. Escala administrativa y MAUP ....... 4:40   (1 + 7)
%   5. Dónde están las fuentes ............ 4:15   (1 + 6)
%   6. La escala de influencia ............ 4:15   (1 + 6)   <- el núcleo
%   7. Conclusiones ....................... 1:00   (1 + 3)
%                                          ------
%                                           23:00
% -----------------------------------------------------------------------------
% VA LARGO: sobran unos 3 minutos, es decir 5 páginas. Candidatas a suprimir,
% en orden de menor coste argumental:
%
%   1. "Funciones de segundo orden"  -> la K residual de la lámina siguiente
%      ya demuestra la agregación; esta solo prepara el terreno.
%   2. "La incertidumbre no es uniforme" -> se dice de viva voz sobre el mapa.
%   3. "Lo que este resultado no dice"   -> la salvedad del 0,41 % de mejora
%      cabe en una frase durante el hallazgo central.
%   4. "Exposición: la media territorial no es la métrica" -> pasa a reserva.
%   5. "Distancia o grafo: el mismo modelo" -> solo si no queda más remedio:
%      es cara de perder porque unifica los bloques 1 y 2 ante este tribunal.
%
% NO SUPRIMIR: "El umbral de declaración sesga la intensidad, no la escala",
% "Esa escala, frente a las unidades en que se legisla", "El hallazgo central"
% y "No es falta de potencia".
% -----------------------------------------------------------------------------
% 5 s: marcan el paso, no se leen.
%
% SI VAS JUSTO DE TIEMPO, suprime en este orden (ninguna sostiene una
% afirmación central):
%   1. "Funciones de segundo orden"   -- control de supuestos
%   2. "Lo que este resultado no dice" -- matiz, se puede decir de viva voz
%   3. "La incertidumbre no es uniforme"
%   4. "Las covariables absorben el 76 %"
% Eso deja el mazo en ~17 min y da margen para preguntas.
%
% LAS CUATRO ÚLTIMAS son de RESERVA (\appendix): no se proyectan, están para
% responder preguntas del tribunal sobre anisotropía, correlación de marcas,
% umbrales de la directiva y residuos del modelo de intensidad.
% =============================================================================
\documentclass[11pt,aspectratio=169]{beamer}

\usepackage[utf8]{inputenc}
\usepackage[T1]{fontenc}
\usepackage[spanish,es-noquoting]{babel}
\usepackage{lmodern}
\usepackage{booktabs}
\usepackage{siunitx}
\usepackage[version=4]{mhchem}
\usepackage{amsmath}
\usepackage{amssymb}
\usepackage{bm}
\usepackage{graphicx}
% Rutas de figuras: LaTeX prueba estas carpetas en orden, así que el mismo
% fichero compila en local (slides/ + ../output/figuras/) y en Overleaf,
% tanto si las imágenes van en output/figuras/ como en figuras/ o en la raíz.
\graphicspath{{../output/figuras/}{output/figuras/}{figuras/}{./}}
\usepackage{tikz}
\usepackage{colortbl}

% =============================================================================
% NOTACIÓN DEL CURSO (Dr. E. A. Chacón, Spatial Statistics, UNI)
% Localizaciones s, desplazamiento d (NO h), proceso latente W y observable Y,
% adyacencia A (la letra W está ocupada por el proceso).
% Centralizado aquí: si cambia una convención, cambia en las 40 páginas.
% =============================================================================
\newcommand{\dom}{D}                         % área de análisis
\newcommand{\loc}{\bm{s}}                    % localización genérica
\newcommand{\loci}[1]{\bm{s}_{#1}}           % localización i-ésima
\newcommand{\dvec}{\vec{d}}                  % desplazamiento (vector)
\newcommand{\dsc}{d}                         % distancia
\newcommand{\Wp}[1]{W(#1)}                   % proceso latente
\newcommand{\Yp}[1]{Y(#1)}                   % variable observable
\newcommand{\varesp}{\sigma^{2}}             % varianza espacial
\newcommand{\pepita}{\tau^{2}}               % pepita / error de medición
\newcommand{\escala}{\phi}                   % parámetro de escala
\newcommand{\suave}{\nu}                     % parámetro de suavizado
\newcommand{\corr}{\rho}                     % función de correlación
\newcommand{\Rphi}{\bm{R}_{\escala}}         % matriz de correlación
\newcommand{\rangop}{r^{*}}                  % rango práctico: corr(r*)=0,05
\newcommand{\Adj}{\bm{A}}                    % matriz de adyacencia
\newcommand{\Prec}{\bm{Q}}                   % matriz de precisión
\newcommand{\Esp}{\mathrm{E}}
\newcommand{\Normal}{\mathcal{N}}

\sisetup{output-decimal-marker={,}, group-separator={\,}}
% Figuras acotadas en ancho y alto. El argumento opcional es la fracción de
% \textheight disponible; bájalo en diapositivas que lleven tabla debajo.
% Esto hace el mazo robusto a la proporción real de cada .png.
\newcommand{\fig}[2][0.52]{%
  \includegraphics[width=\linewidth,height=#1\textheight,keepaspectratio]{#2}}

\newcommand{\ugm}{\si{\micro\gram\per\cubic\metre}}

% =============================================================================
% Tema formal propio: sin adornos, cintillo de sección y pie institucional
% =============================================================================
\usetheme{default}
\usecolortheme{default}
\setbeamertemplate{navigation symbols}{}

\definecolor{azul}{RGB}{28,72,122}
\definecolor{azulc}{RGB}{112,148,190}
\definecolor{rojo}{RGB}{160,25,45}
\definecolor{gris}{RGB}{105,105,105}
\definecolor{grisc}{RGB}{240,240,243}

\setbeamercolor{structure}{fg=azul}
\setbeamercolor{frametitle}{fg=azul}
\setbeamercolor{title}{fg=azul}
\setbeamercolor{block title}{bg=grisc,fg=azul}
\setbeamercolor{block body}{bg=grisc,fg=black}
\setbeamercolor{block title alerted}{bg=grisc,fg=rojo}
\setbeamercolor{block body alerted}{bg=grisc,fg=black}
\setbeamercolor{alerted text}{fg=rojo}
\setbeamerfont{frametitle}{series=\bfseries,size=\large}
\setbeamerfont{framesubtitle}{size=\small,shape=\itshape}
\setbeamercolor{framesubtitle}{fg=gris}
\setbeamertemplate{itemize item}{\color{azul}\rule[0.35ex]{0.5ex}{0.5ex}}
\setbeamertemplate{itemize subitem}{\color{azulc}\rule[0.35ex]{0.4ex}{0.4ex}}
\setbeamertemplate{blocks}[rounded][shadow=false]

% Cintillo superior con el nombre de la sección
\setbeamertemplate{headline}{%
  \vspace{2.2mm}%
  \hbox{\hspace*{5mm}{\scriptsize\color{gris}\insertsectionhead}\hfill}%
  \vspace{1mm}%
  \hbox{\hspace*{5mm}{\color{azulc}\rule{0.95\paperwidth}{0.4pt}}}%
}

% Pie institucional
\setbeamertemplate{footline}{%
  \hbox{\hspace*{5mm}{\color{azulc}\rule{0.95\paperwidth}{0.4pt}}}%
  \vspace{1mm}%
  \hbox{\hspace*{5mm}{\tiny\color{gris}R. Fernández \textperiodcentered\
  J. Carhuaricra\quad\textbar\quad
  \ce{NO2} en España: la escala de la influencia industrial}%
  \hfill{\tiny\color{gris}\insertframenumber}\hspace*{5mm}}%
  \vspace{2mm}%
}

% Divisoria de sección
\AtBeginSection[]{%
  \begin{frame}[plain]
    \vfill
    \hspace*{12mm}\begin{minipage}{0.78\textwidth}
      {\color{azulc}\rule{\textwidth}{1pt}}\\[3mm]
      {\small\color{gris}Parte \insertsectionnumber}\\[2mm]
      {\LARGE\color{azul}\bfseries\insertsectionhead}\\[3mm]
      {\color{azulc}\rule{\textwidth}{1pt}}
    \end{minipage}
    \vfill
  \end{frame}%
}

\begin{document}

% =============================================================================
\begin{frame}[plain]
  \vfill
  \hspace*{8mm}\begin{minipage}{0.85\textwidth}
    {\color{azulc}\rule{\textwidth}{1.2pt}}\\[5mm]
    {\LARGE\color{azul}\bfseries La escala importa}\\[3mm]
    {\large Convergencia entre el agrupamiento de fuentes industriales\\
     y su influencia sobre el \ce{NO2} de fondo en España}\\[8mm]
    {\color{azulc}\rule{\textwidth}{0.5pt}}\\[4mm]
    \begin{minipage}[t]{0.48\textwidth}
      Randy Fernández Palomino\\[0.5mm]
      {\footnotesize\color{gris}\texttt{randy.fernandez.p@uni.pe}}
    \end{minipage}%
    \begin{minipage}[t]{0.48\textwidth}
      José Carhuaricra Cusipuma\\[0.5mm]
      {\footnotesize\color{gris}\texttt{jose.carhuaricra@unh.edu.pe}}
    \end{minipage}\\[6mm]
    {\small\color{gris}Estadística Espacial \textperiodcentered\ Doctorado en
    Ciencias e Ingeniería Estadística}\\[1mm]
    {\small\color{gris}Universidad Nacional de Ingeniería (UNI)
    \hfill 2026}
  \end{minipage}
  \vfill
\end{frame}

% =============================================================================
\begin{frame}{Contenido}
\vspace{-5mm}
\begin{columns}[T]
\begin{column}{0.46\textwidth}
\small
\begin{enumerate}
  \item El problema
  \item Datos y diseño
  \item El campo de fondo
  \item Dónde están las fuentes
\end{enumerate}
\end{column}
\begin{column}{0.46\textwidth}
\small
\begin{enumerate}
  \setcounter{enumi}{4}
  \item La escala de influencia
  \item Escala administrativa y MAUP
  \item Conclusiones
\end{enumerate}
\end{column}
\end{columns}

\vspace{3mm}
\centering
\scriptsize
\begin{tabular}{@{}llc@{}}
\toprule
\textbf{Rama} & \textbf{Conjunto índice} & \textbf{Objeto} \\
\midrule
Variación espacial continua & $\loc \in \dom \subset \mathbb{R}^2$, continuo
  & $\Yp{\loc}$ \\
Variación espacial discreta & $\{R_1,\dots,R_{50}\}$, retícula irregular
  & $\Yp{R_i}$ \\
Procesos puntuales & el índice \emph{es} la aleatoriedad & $\{\loci{i}\}$ \\
\bottomrule
\end{tabular}

\vspace{3mm}
\begin{block}{Hilo conductor}
\footnotesize
Lo que cambia entre las tres ramas es \textbf{el conjunto índice}. Las articula
una pregunta que ninguna responde sola: \textbf{¿a qué escala espacial opera la
influencia industrial sobre el \ce{NO2} de fondo?}
\end{block}
\end{frame}

% #############################################################################
\section{El problema}
% #############################################################################

\begin{frame}{El problema está mal condicionado}
\vspace{-5mm}
\begin{columns}[T]
\begin{column}{0.54\textwidth}
El \ce{NO2} tiene \textbf{dos firmas simultáneas}:
\begin{itemize}
  \item emisiones puntuales industriales
  \item emisiones lineales de tráfico
\end{itemize}

\vspace{4mm}
Los modelos de referencia (ESCAPE, LUR europeos) alcanzan
$R^2 \approx \numrange{0.5}{0.8}$\dots

\vspace{2mm}
\textbf{\dots pero no separan mecanismos.}

\vspace{5mm}
{\footnotesize\color{gris}En la llanura padana, la cuenca del Ruhr o la Alta
Silesia, industria y ciudad son la misma covariable.}
\end{column}

\begin{column}{0.42\textwidth}
\vspace{-4mm}
\begin{block}{No es un problema de método}
Es de \textbf{identificabilidad}.

\vspace{2mm}
Ningún estimador puede separar dos covariables casi colineales, por
sofisticado que sea.
\end{block}

\vspace{4mm}
{\footnotesize\color{gris}Beelen et al.\ (2013); Chen et al.\ (2019);
Vizcaíno y Lavalle (2018).}
\end{column}
\end{columns}

\end{frame}

% =============================================================================
\begin{frame}{España como diseño natural}

\begin{columns}[T]
\begin{column}{0.5\textwidth}
\fig[0.92]{20_patron_industrial.png}
\end{column}

\begin{column}{0.46\textwidth}
\vspace{1mm}
{\small Los grandes emisores están lejos de los grandes focos urbanos.}

\vspace{3mm}
{\small\textbf{Industriales:} Huelva, Tarragona, Puertollano, As Pontes,
Avilés, Algeciras.}

\vspace{2mm}
{\small\textbf{Urbanos:} Madrid, Barcelona.}

\vspace{3mm}
\begin{block}{}
\small
Un país donde el efecto industrial \textbf{puede} estimarse sin confundirse
con la urbanización.
\end{block}

\vspace{2mm}
{\scriptsize\color{gris}Canarias concentra el \SI{24}{\percent} del
\ce{NOx} industrial y queda fuera del dominio: las tres ramas exigen
conexidad.}
\end{column}
\end{columns}

\end{frame}

% =============================================================================
\begin{frame}{La pregunta y el diseño de contraste}

\begin{center}
\large
¿A qué \textbf{escala espacial} opera la influencia industrial?
\end{center}

\vspace{4mm}
\begin{columns}[T]
\begin{column}{0.46\textwidth}
\begin{block}{Rama B --- proceso puntual}
Observa \textbf{solo coordenadas} de \num{297} instalaciones.

\vspace{2mm}
\emph{¿A qué distancia se agrupan entre sí?}
\end{block}
\end{column}
\begin{column}{0.46\textwidth}
\begin{block}{Rama C --- acoplamiento}
Observa \textbf{concentraciones} en \num{212} estaciones de fondo.

\vspace{2mm}
\emph{¿A qué distancia deja de notarse la emisión?}
\end{block}
\end{column}
\end{columns}

\vspace{4mm}
\begin{center}
\begin{tikzpicture}
\node[draw=rojo,line width=0.7pt,rounded corners=2pt,inner sep=3mm,
      fill=grisc]{%
\begin{minipage}{0.74\textwidth}\centering
{\color{rojo}\textbf{No hay ninguna razón aritmética para que coincidan.}}\\[1mm]
{\small Si convergen, la unidad de influencia es el distrito, no la
instalación.}
\end{minipage}};
\end{tikzpicture}
\end{center}

\end{frame}

% #############################################################################
\section{Datos y diseño}
% #############################################################################

\begin{frame}{Siete fuentes, todas públicas y de acceso programático}

\vspace{-5mm}
\centering
\small
\begin{tabular}{@{}lll@{}}
\toprule
\textbf{Variable} & \textbf{Fuente} & \textbf{Resolución} \\
\midrule
Geometrías \textsc{nuts} 2021 & EEA & \num{59} \textsc{nuts}-3 \\
\ce{NO2} por estación (2023) & EEA, dataset E1a & \num{462} estaciones \\
Clasificación de estaciones & EEA DiscoData & --- \\
Instalaciones IED y emisiones & EEA DiscoData & \num{297} sitios \\
\ce{NO2} troposférico & Copernicus Sentinel-5P & $\sim$\SI{5}{\kilo\metre} \\
Cobertura del suelo & CORINE CLC2018 & \SI{100}{\metre} \\
Densidad de población & Eurostat & \textsc{nuts}-3 \\
\bottomrule
\end{tabular}

\vspace{7mm}
\begin{columns}
\begin{column}{0.46\textwidth}
\centering
{\footnotesize Dominio: península y Baleares\\
\SI{498522}{\kilo\metre\squared} \textperiodcentered\ \textsc{epsg}:3035}
\end{column}
\begin{column}{0.46\textwidth}
\centering
{\footnotesize Licencias: CC-BY 4.0 (EEA, Eurostat)\\
aviso legal Copernicus (Sentinel, CORINE)}
\end{column}
\end{columns}

\end{frame}

% =============================================================================
\begin{frame}{La red de medición está sesgada}
\vspace{-10mm}
\begin{columns}[T]
\begin{column}{0.5\textwidth}
\fig[1.5]{01_histograma.png}
\end{column}
\begin{column}{0.46\textwidth}
\vspace{3mm}
\textbf{Desagrupamiento por celdas de \SI{25}{\kilo\metre}}

\vspace{3mm}
\centering
\begin{tabular}{@{}lc@{}}
\toprule
Media bruta & \SI{11.36}{\ugm} \\
\rowcolor{grisc}
Media desagrupada & \textbf{\SI{8.58}{\ugm}} \\
\bottomrule
\end{tabular}

\vspace{4mm}
\raggedright
{\small Sin corregir, la concentración de estaciones en zonas urbanas sesga al
alza toda la superficie predicha.}

\vspace{3mm}
{\footnotesize\color{gris}Variable de trabajo: $\log(\ce{NO2})$. Asimetría
bruta \num{0.81}; tras el logaritmo, \num{-0.58}.}
\end{column}
\end{columns}

\end{frame}

% =============================================================================
\begin{frame}{Dos decisiones que condicionan todo lo demás}
\vspace{-5mm}
\begin{block}{1. Solo estaciones de fondo alimentan el modelo geoestadístico}
\small
\num{212} de \num{462}. Las de tráfico miden a microescala: el \ce{NO2} cae
\SIrange{60}{80}{\percent} en los primeros \SI{50}{\metre} desde la calzada.
No son realizaciones del campo regional.

\vspace{1.5mm}
$\rightarrow$ Se reservan como \textbf{validación externa}. El sesgo medido en
ellas \emph{cuantifica el recargo local}: deja de ser un problema y pasa a ser
un resultado.
\end{block}

\vspace{3mm}
\begin{block}{2. La respuesta areal es la media de bloque, no la exposición ponderada}
\small
La exposición ponderada por población correlaciona $\rho = \num{0.948}$ con la
densidad \textbf{por construcción}: se calcula con la misma distribución que
después entra como predictor.

\vspace{1.5mm}
$\rightarrow$ Regresarla sobre densidad sería circular. Se conserva como
métrica de política pública, no como variable dependiente.
\end{block}

\end{frame}

% #############################################################################
\section{El campo de fondo}
% #############################################################################

\begin{frame}{Estructura espacial: el variograma}
\vspace{-10mm}
\begin{columns}[T]
\begin{column}{0.5\textwidth}
\vspace{4mm}
\fig[0.94]{04_variograma_ajustado.png}
\end{column}
\begin{column}{0.46\textwidth}
\vspace{6mm}
\centering
\begin{tabular}{@{}lc@{}}
\toprule
Rango práctico $\rangop$ & \SI{72}{\kilo\metre} \\
Pepita $\pepita/(\varesp+\pepita)$ & \SI{9.6}{\percent} \\
\bottomrule
\end{tabular}

\vspace{3mm}
\raggedright
{\footnotesize Rango práctico en el sentido $\corr(\rangop)=\num{0.05}$. El
\SI{90.4}{\percent} de la variabilidad está espacialmente estructurada.}

\vspace{2mm}
{\footnotesize\color{gris}\textbf{Advertencia declarada:} el ajuste es el
límite suave de la familia de Matérn ($\suave$ grande) e implica variación por
debajo de la menor distancia entre estaciones. Poco creíble en calidad del
aire: se declara y se controla con el variograma residual.}
\end{column}
\end{columns}

\end{frame}

% =============================================================================
\begin{frame}{Las covariables absorben el \SI{76}{\percent} de la varianza}

\begin{columns}[T]
\begin{column}{0.48\textwidth}
\centering
{\footnotesize Variograma omnidireccional}\\[1mm]
\fig[0.48]{02_variograma_omni.png}
\end{column}
\begin{column}{0.48\textwidth}
\centering
{\footnotesize Variograma de los residuos tras la deriva}\\[1mm]
\fig[0.48]{05_variograma_residual.png}
\end{column}
\end{columns}

\vspace{-1mm}
\begin{center}
$\Yp{\loci{i}} = \underbrace{\bm{x}^{\!\top}\!(\loci{i})\bm{\beta}}
_{\text{\footnotesize variación espacial explicada}}
+\;\underbrace{\Wp{\loci{i}}}
_{\text{\footnotesize no explicada}}\; +\; \varepsilon_i$
\end{center}

\vspace{-1mm}
\begin{center}
Rango bruto \SI{72}{\kilo\metre} $\;\longrightarrow\;$ rango residual
\textbf{\SI{26}{\kilo\metre}}
\end{center}
\vspace{-2mm}
\begin{center}
{\footnotesize Varianza bruta \num{0.650} $\rightarrow$ residual \num{0.155}.
Una tendencia infla el semivariograma con un término cuadrático en la
distancia: por eso el rango se estima sobre residuos.}
\end{center}

\end{frame}

% =============================================================================
\begin{frame}{Kriging ordinario frente a kriging con deriva externa}
\vspace{-2mm}
\begin{columns}[T]
\begin{column}{0.48\textwidth}
\centering
{\footnotesize Kriging ordinario}\\[1mm]
\fig[0.45]{06_mapa_ko.png}
\end{column}
\begin{column}{0.48\textwidth}
\centering
{\footnotesize Kriging con deriva externa}\\[1mm]
\fig[0.45]{07_mapa_kde.png}
\end{column}
\end{columns}

\vspace{-3mm}
\begin{center}
\footnotesize
\begin{tabular}{@{}lccc@{}}
\toprule
 & \textbf{RMSE} & \textbf{$R^2$} & \textbf{Sesgo} \\
\midrule
Kriging ordinario & \num{5.565} & \num{0.481} & \num{0.933} \\
\rowcolor{grisc}
Con deriva externa & \textbf{\num{3.657}} & \textbf{\num{0.776}} &
\textbf{\num{0.156}} \\
\midrule
Mejora & \SI{-34.3}{\percent} & & \\
\bottomrule
\end{tabular}
\end{center}

% ---------------------------------------------------------------------------
% ---------------------------------------------------------------------------
\vspace{-2mm}
\centering{\footnotesize Validación \emph{leave-one-out}, \num{212}
estaciones. Cobertura al \SI{95}{\percent}: \num{0.910}, con
$\mathrm{ESD}=\num{1.225}$: la varianza de kriging es un \SI{50}{\percent}
menor de la real, porque trata $\hat\varesp$, $\hat\pepita$ y $\hat\escala$
como conocidos.}

\end{frame}

% =============================================================================
\begin{frame}{La incertidumbre no es uniforme}
\vspace{-5mm}
\begin{columns}[T]
\begin{column}{0.52\textwidth}
\fig[1.2]{08_mapa_error.png}
\end{column}
\begin{column}{0.44\textwidth}
\vspace{6mm}
{\small El error de predicción crece donde la red de medición es rala:
interior de Castilla, Extremadura, Aragón.}

\vspace{4mm}
{\small Toda afirmación sobre esas zonas debe leerse con el intervalo, no con
la media puntual.}

\vspace{6mm}
\begin{block}{}
\small
Propagar la incertidumbre es lo que distingue un mapa de una interpolación
decorativa.
\end{block}
\end{column}
\end{columns}

\end{frame}

% =============================================================================
\begin{frame}{El resultado que importa: el recargo local}
\framesubtitle{Medido en las \num{250} estaciones excluidas del ajuste}

\vspace{1mm}
\centering
\small
\begin{tabular}{@{}llcc@{}}
\toprule
\textbf{Tipo} & \textbf{Entorno} & \textbf{$n$} & \textbf{Recargo (\ugm)} \\
\midrule
Industrial & rural & \num{32} & \num{0.58} \\
Industrial & suburbano & \num{61} & \num{0.76} \\
Industrial & urbano & \num{37} & \num{1.52} \\
\midrule
\rowcolor{grisc}
Tráfico & suburbano & \num{18} & \num{4.77} \\
\rowcolor{grisc}
Tráfico & urbano & \num{102} & \num{3.00} \\
\bottomrule
\end{tabular}

\vspace{6mm}
\begin{block}{}
\centering
\small
En entorno rural las estaciones industriales miden un \SI{52}{\percent} más
que las de fondo (\num{5.95} frente a \num{3.91}\,\ugm). En entorno urbano la
diferencia se anula.

\vspace{2mm}
\textbf{La industria eleva el \ce{NO2} donde el fondo es bajo, pero queda
sepultada por el tráfico en las ciudades.}
\end{block}

\end{frame}

% #############################################################################
\section{Escala administrativa y MAUP}
% MAUP = Modifiable Areal Unit Problem (problema de la unidad de área
% modificable): los resultados estadísticos dependen del tamaño y la forma
% de las unidades de agregación elegidas.
% #############################################################################

\begin{frame}{Distancia o grafo: el mismo modelo}

\vspace{2mm}
\begin{columns}[T]
\begin{column}{0.47\textwidth}
\centering
\textbf{Rama continua}
\[ \bm{Y} \sim \Normal_n\!\left(\bm{X}\bm{\beta},\;
   \varesp\Rphi + \pepita\bm{I}_n\right) \]
\vspace{1mm}
{\small La covarianza se construye con \textbf{distancias}: $\Rphi$ depende de
$\lVert\loci{i}-\loci{j}\rVert$.}
\end{column}
\begin{column}{0.47\textwidth}
\centering
\textbf{Rama areal}
\[ \bm{Y}\mid\bm{X} \sim \Normal\!\left(\bm{X}\bm{\beta},\;
   \Prec^{-1} + \varesp\bm{I}\right) \]
\vspace{1mm}
{\small La covarianza se construye con un \textbf{grafo}:
$Q_{ij}\neq 0 \iff$ $R_i$ y $R_j$ son vecinas.}
\end{column}
\end{columns}

\vspace{4mm}
\begin{block}{}
\centering
\small
Misma estructura: $\Normal(\bm{X}\bm{\beta},\ \text{espacial}+\text{ruido})$.
Solo cambia cómo se especifica el efecto espacial.

\vspace{1.5mm}
La covarianza es \textbf{densa} e informa de la dependencia marginal; la
precisión es \textbf{dispersa} e informa de la independencia condicional.
\end{block}

\vspace{2mm}
\centering{\footnotesize\color{gris}No se modela dos veces lo mismo: se modela
el mismo campo bajo dos conjuntos índice, y la comparación entre ambos es el
objeto de esta parte.}

\end{frame}

% =============================================================================
\begin{frame}{Dependencia espacial: en los residuos, no en la respuesta}
\vspace{-2mm}
\begin{columns}[T]
\begin{column}{0.48\textwidth}
\centering
{\footnotesize Diagrama de dispersión de Moran}\\[-0.5mm]
\fig[0.45]{40_moran_scatter.png}
\end{column}
\begin{column}{0.48\textwidth}
\centering
{\footnotesize Agrupamientos locales (\textsc{lisa})}\\[-0.5mm]
\fig[0.45]{41_lisa.png}
\end{column}
\end{columns}

\vspace{-3mm}
\begin{center}
\footnotesize
\begin{tabular}{@{}lccc@{}}
\toprule
 & \textbf{$I$ de Moran} & \textbf{$\Esp[I]$ bajo $H_0$} & \textbf{$p$} \\
\midrule
Sobre la respuesta (media areal) & \num{0.013} & \num{-0.020} & \num{0.32} \\
\rowcolor{grisc}
Sobre los residuos del \textsc{mco} & \num{0.187} & \num{-0.020} & \num{0.0019} \\
\bottomrule
\end{tabular}
\end{center}

\vspace{-1mm}
\centering{\small Con $n=\num{50}$, $\Esp[I]=-1/(n-1)=\num{-0.020}$: el
valor sobre la respuesta no es «cero», es su esperanza bajo independencia. La
dependencia aparece \textbf{después} de condicionar, no antes: eso justifica el
modelo de error espacial, no un contagio entre provincias.
{\scriptsize\color{gris} \textsc{lisa}: indicadores locales de asociación
espacial, la descomposición local del índice de Moran.}}

\end{frame}

% =============================================================================
\begin{frame}{A escala provincial, la industria desaparece}
\framesubtitle{Cuatro intentos de refutación de un resultado nulo}

\vspace{1mm}
\centering
\small
\begin{tabular}{@{}lcc@{}}
\toprule
\textbf{Especificación} & \textbf{$\hat\beta$ industria} & \textbf{$p$} \\
\midrule
\textsc{mco}, emisiones totales & \num{-0.019} & \num{0.63} \\
\textsc{mco}, presión a \SI{10}{\kilo\metre} & \num{-0.043} & \num{0.36} \\
\textsc{mco} sin Madrid (Cook $=\num{0.97}$) & \num{0.004} & \num{0.90} \\
\textsc{sar} con regresión (error espacial) & \num{0.017} & \num{0.68} \\
\bottomrule
\end{tabular}

\vspace{3mm}
\begin{columns}[T]
\begin{column}{0.46\textwidth}
\begin{block}{\textsc{sar} con regresión}
\centering
\small
$\bm{U}=\bm{B}\bm{U}+\bm{\varepsilon}$, \; $\bm{U}=\bm{Y}-\bm{X}\bm{\beta}$

\vspace{1mm}
$\bm{B}=\alpha\,\bm{D}_a^{-1}\Adj$, \;
$\hat\alpha=\num{0.446}\in(-1,1)$
\end{block}
\end{column}
\begin{column}{0.46\textwidth}
\begin{block}{Densidad de población}
\centering
$\hat\beta = \num{0.575}$, \quad $p = \num{2.1e-9}$

\vspace{1mm}
{\footnotesize Todos los \textsc{vif} $< \num{1.8}$}
\end{block}
\end{column}
\end{columns}

\vspace{3mm}
\centering{\scriptsize \textsc{mco}: mínimos cuadrados ordinarios.
\textsc{sem}: \emph{spatial error model}, el \textsc{sar} con regresión sobre
el residuo. Se reporta \textsc{sem} y no el modelo Durbin pese a su \textsc{aic}
marginalmente inferior ($\Delta$AIC $= \num{0.25}$): los contrastes robustos
separan hacia el error, el $\rho$ del retardo no es significativo
($p = \num{0.91}$) y ningún impacto del Durbin alcanza significación.}

\end{frame}

% =============================================================================
\begin{frame}{MAUP: los dos componentes}
\framesubtitle{\emph{Modifiable areal unit problem}}

\begin{columns}[T]
\begin{column}{0.5\textwidth}
\fig[0.7]{43_maup_zonificacion.png}
\end{column}
\begin{column}{0.46\textwidth}
{\footnotesize\textbf{Escala.} A \textsc{nuts}-2 aparece
$\hat\beta = \num{-0.562}$ ($p = \num{0.006}$), pero es íntegramente Madrid:
Cook $=\num{4.64}$, diecinueve veces el umbral. Sin ella,
$\hat\beta = \num{-0.063}$ ($p = \num{0.68}$) y el $R^2$ ajustado cae de
\num{0.775} a \num{0.597}, casi el \num{0.601} de \textsc{nuts}-3.}

\vspace{2mm}
{\footnotesize\textbf{Zonificación.} \num{1000} particiones de \num{16}
grupos por $k$-medias:}

\vspace{1.5mm}
\centering
\footnotesize
\begin{tabular}{@{}lc@{}}
\toprule
Densidad significativa en & \SI{100}{\percent} \\
\rowcolor{grisc}
Industria significativa en & \SI{2.5}{\percent} \\
\bottomrule
\end{tabular}

\vspace{2mm}
\raggedright
{\scriptsize El valor de \textsc{nuts}-2 queda \textbf{fuera del rango
completo} de las particiones. Madrid \textbf{nunca} queda aislada en las mil
réplicas: \textsc{nuts}-2 la aísla por razones históricas, no geométricas.}
\end{column}
\end{columns}

\end{frame}

% =============================================================================
\begin{frame}{Los mecanismos son inseparables a esta escala}
\vspace{-5mm}
\begin{columns}[T]
\begin{column}{0.5\textwidth}
\fig[0.8]{51_sintesis_bloques.png}
\end{column}
\begin{column}{0.46\textwidth}
\vspace{2mm}
\centering
\small
\begin{tabular}{@{}lcc@{}}
\toprule
\textbf{Bloque} & \textbf{único} & \textbf{total} \\
\midrule
Uso del suelo & \num{0.157} & \num{0.906} \\
Urbanización & \num{0.023} & \num{0.664} \\
Industria & \num{0.005} & \num{0.287} \\
Geografía & \num{0.001} & \num{0.120} \\
\bottomrule
\end{tabular}

\vspace{3mm}
\begin{alertblock}{}
\centering
Varianza única total:
\textbf{\SI{19.7}{\percent}} del $R^2$
\end{alertblock}

\vspace{1mm}
{\footnotesize El \SI{80.3}{\percent} restante es \textbf{compartida}: no
atribuible a ningún mecanismo por separado.}
\end{column}
\end{columns}

\vspace{1mm}
\centering{\footnotesize La industria se ubica donde hay ciudades, puertos y
autopistas: las mismas variables que generan \ce{NO2} de tráfico.
\textbf{La endogeneidad es una propiedad del territorio, no un defecto del
análisis.}}

\end{frame}

% =============================================================================
\begin{frame}{Exposición: la media territorial no es la métrica}
\vspace{-1mm}
\begin{columns}[T]
\begin{column}{0.48\textwidth}
\centering
{\footnotesize Exposición ponderada por población}\\[1mm]
\fig[0.6]{13_exposicion_poblacional.png}
\end{column}
\begin{column}{0.48\textwidth}
\centering
{\footnotesize Brecha frente a la media areal}\\[1mm]
\fig[0.6]{15_brecha_exposicion.png}
\end{column}
\end{columns}

\vspace{-2.5mm}
\centering
\small
\begin{tabular}{@{}lccc@{}}
\toprule
 & \textbf{areal} & \textbf{ponderada} &
 \textbf{\% pobl.\ $>$ \SI{20}{\ugm}} \\
\midrule
Madrid & \num{7.52} & \num{15.01} & \SI{37.3}{\percent} \\
Barcelona & \num{6.64} & \num{14.21} & \SI{30.6}{\percent} \\
Bizkaia & \num{5.74} & \num{11.80} & \SI{22.7}{\percent} \\
\bottomrule
\end{tabular}

\end{frame}

% =============================================================================
\begin{frame}{Cuánto, no solo si}

\vspace{2mm}
\centering
\small
\begin{tabular}{@{}lc@{}}
\toprule
 & \textbf{\ugm} \\
\midrule
\rowcolor{grisc}
\textbf{Cota superior del efecto industrial areal} & \textbf{\num{0.96}} \\
Recargo local en estaciones industriales & $\approx\num{1.0}$ \\
Recargo local en estaciones de tráfico & $\approx\num{3.5}$ \\
Media areal nacional de fondo & \num{4.28} \\
\bottomrule
\end{tabular}

\vspace{4mm}
\begin{block}{Por qué una cota y no un $p$-valor}
\centering
No es que \emph{ignoremos} si la industria influye.

\vspace{1mm}
Sabemos que \textbf{no puede influir más que eso}: el \SI{22.3}{\percent} de
la media nacional.
\end{block}

\vspace{3mm}
\centering{\footnotesize El \SI{43.3}{\percent} de las estaciones de tráfico superan
hoy los \SI{20}{\ugm}, frente al \SI{0.16}{\percent} de la superficie de
fondo. España cumple con holgura la directiva vigente, pero el problema reside
en una escala por debajo de la que este trabajo puede resolver.}

\end{frame}

% #############################################################################
\section{Dónde están las fuentes}
% #############################################################################

\begin{frame}{El patrón industrial no es aleatorio}

\begin{columns}[T]
\begin{column}{0.5\textwidth}

\fig[1.9]{22_cuadrantes.png}
\end{column}
\begin{column}{0.46\textwidth}
{\scriptsize \num{297} instalaciones sobre \SI{504327}{\kilo\metre\squared}.}

\vspace{2mm}
\centering
\small
\begin{tabular}{@{}lc@{}}
\toprule
Intensidad media & \num{0.589}/\SI{1000}{\kilo\metre\squared} \\
Distancia al vecino & \SI{12.7}{\kilo\metre} \\
Esperada bajo \textsc{ace}$^{\dagger}$ & \SI{20.6}{\kilo\metre} \\
\bottomrule
\end{tabular}

\vspace{3mm}
\begin{alertblock}{Contraste de cuadrantes}
\centering
$X^2 = \textstyle\sum_{i} (n_i-e_i)^2/e_i = \num{531.64}$

\vspace{0.5mm}
$p = \num{0.002}$

\vspace{0.5mm}
{\scriptsize Monte Carlo, $m=\num{52}$ teselas. $H_0$: proceso homogéneo.}
\end{alertblock}

\vspace{1mm}
{\scriptsize\color{gris}$^{\dagger}$\,\textsc{ace}: aleatoriedad espacial
completa, o sea un Poisson homogéneo. Deduplicado por \texttt{InspireSiteId}:
la teoría exige un proceso \emph{ordenado}.}
\end{column}
\end{columns}

\end{frame}

% =============================================================================
\begin{frame}{Intensidad modelada: qué predice la ubicación}

\vspace{-1mm}
\begin{center}
$\log \lambda(\loc) = B(\loc) + \theta_1 Z_1(\loc) + \cdots +
\theta_p Z_p(\loc)$
\end{center}

\vspace{1mm}
\begin{columns}[T]
\begin{column}{0.45\textwidth}
\fig[1.90]{21_intensidad.png}
\end{column}
\begin{column}{0.5\textwidth}
\centering
\small
\begin{tabular}{@{}lcc@{}}
\toprule
\textbf{$Z_j(\loc)$, efecto por rango intercuartílico} & \textbf{Factor} & $p$ \\
\midrule
Suelo industrial & $\times\num{3.94}$ & $<\num{0.001}$ \\
Superficie artificial & $\times\num{1.32}$ & $<\num{0.001}$ \\
Transporte & $\times\num{0.75}$ & $<\num{0.001}$ \\
\rowcolor{yellow!30}
\textbf{Densidad de población} & $\times\num{0.999}$ &
\textbf{\num{0.997}} \\
\bottomrule
\end{tabular}

\vspace{2mm}
\begin{block}{}
\scriptsize
Condicionando en uso del suelo, \textbf{dónde vive la gente no predice dónde
se instala la industria}: el factor \num{0.999} es el valor de referencia bajo
independencia. Evidencia directa contra la colinealidad total entre ambos
mecanismos.
\end{block}

\vspace{1mm}
{\scriptsize\color{gris}Los $Z_j(\loc)$ han de conocerse en \textbf{todo} el
dominio: de ahí que sean rásteres y no atributos de planta.}
\end{column}
\end{columns}

\end{frame}

% =============================================================================
\begin{frame}{Funciones de segundo orden: hay agregación}

\begin{columns}[T]
\begin{column}{0.48\textwidth}
\centering
{\footnotesize $K$ inhomogénea}\\[-0.5mm]
\fig[0.58]{23_Kinhom.png}
\end{column}
\begin{column}{0.48\textwidth}
\centering
{\footnotesize Correlación por pares $g$}\\[-0.5mm]
\fig[0.58]{24_ginhom.png}
\end{column}
\end{columns}

\vspace{-2mm}
\begin{center}
$g_2(\loc,\loc') = \dfrac{\lambda_2(\loc,\loc')}{\lambda(\loc)\lambda(\loc')}
\;\;\overset{\text{independencia}}{=}\;\; 1,
\qquad
K(t) = \pi t^2 \;\;\text{bajo Poisson homogéneo}$
\end{center}

\vspace{-1mm}
\begin{center}
{\footnotesize Rechazada la homogeneidad, la caracterización de segundo orden
ha de condicionar en $\hat\lambda(\loc)$: de ahí las versiones
\textbf{inhomogéneas}. Envolventes con \num{199} simulaciones; la curva
empírica sale de la envolvente.}
\end{center}

\end{frame}

% =============================================================================
\begin{frame}{Tras condicionar en todo, quedan distritos}
\vspace{-10mm}
\begin{columns}[T]
\begin{column}{0.5\textwidth}
\fig[1.5]{26_Kinhom_residual.png}
\end{column}
\begin{column}{0.46\textwidth}
\vspace{3mm}
{\small La $K$ inhomogénea sobre la intensidad \emph{ajustada} sigue mostrando
agregación no explicada hasta \textbf{\SI{38}{\kilo\metre}}.}

\vspace{3mm}
{\small Las covariables observadas —uso del suelo, transporte, población— no
capturan toda la estructura.}

\vspace{5mm}
\begin{block}{}
\small
Interpretación: infraestructura compartida, cadenas de suministro, herencia
histórica.

\vspace{1.5mm}
\textbf{No es ruido: es estructura.}
\end{block}
\end{column}
\end{columns}

\end{frame}

% =============================================================================
\begin{frame}{Proceso agregado: padres no observados y descendencia}

\vspace{-2mm}
\begin{center}
{\footnotesize Padres de un Poisson de intensidad $\kappa$ $\rightarrow$
descendencia por padre $\rightarrow$ \textbf{lo observado son los hijos}. Con
dispersión gaussiana, es el proceso de Thomas.}
\end{center}

\begin{center}
\small
\begin{tabular}{@{}lc@{}}
\toprule
Escala del conglomerado $\sigma$ & \SI{2.7}{\kilo\metre} \\
\rowcolor{grisc}
\textbf{Diámetro efectivo del distrito} $(4\sigma)$ &
\textbf{\SI{10.7}{\kilo\metre}} \\
Conglomerados estimados en España & \num{520} \\
Probabilidad de hermandad $p_{\text{sib}}$ & \num{0.915} \\
Fuerza del conglomerado $\varphi$ & \num{10.74} \\
\bottomrule
\end{tabular}
\end{center}

\vspace{1mm}
\begin{center}
\begin{tikzpicture}[scale=0.80]
  \foreach \x/\y in {0/0, 3.6/0.35, 7.2/-0.25} {
    \draw[azulc,dashed] (\x,\y) circle (1.05);
    \fill[azul] (\x,\y) circle (0.055);
    \foreach \dx/\dy in {0.38/0.28, -0.45/0.18, 0.18/-0.45, -0.28/-0.35,
                         0.55/-0.08} {
      \fill[rojo] (\x+\dx,\y+\dy) circle (0.05);}
  }
  \node[gris,anchor=west] at (8.6,0) {\footnotesize $\varnothing \approx
  \SI{10.7}{\kilo\metre}$};
\end{tikzpicture}
\end{center}

\centering{\scriptsize\color{gris}$\sigma$ es la desviación típica de la
dispersión en torno a un centro latente; $4\sigma$ contiene $\approx$ el
\SI{95}{\percent}. Un proceso agregado es un Poisson \textbf{doblemente
estocástico}: su intensidad $\Lambda(\loc)$ es aleatoria y no se observa.}

\end{frame}

% =============================================================================
\begin{frame}{El umbral de declaración sesga la intensidad, no la escala}

\vspace{2mm}
\begin{columns}[T]
\begin{column}{0.5\textwidth}
\begin{block}{El problema}
\small
El registro E-PRTR solo recoge instalaciones que declaran más de
\SI{100}{\tonne} anuales de \ce{NOx}. El patrón observado $Y$ es un
\textbf{subconjunto} del patrón industrial real $Z$.
\end{block}
\end{column}
\begin{column}{0.46\textwidth}
\begin{block}{El resultado}
\small
Si cada evento de $Z$ se retiene de forma independiente con probabilidad $p$:
\[ \lambda_Y(\loc) = p\,\lambda_Z(\loc) \qquad\text{pero}\qquad
   K_Y(t) = K_Z(t) \]
\end{block}
\end{column}
\end{columns}

\vspace{3mm}
\begin{center}
\begin{tikzpicture}
\node[draw=azul,line width=0.7pt,rounded corners=2pt,inner sep=2.5mm,
      fill=grisc]{%
\begin{minipage}{0.80\textwidth}\centering
{\color{azul}El hallazgo central depende de $K$, no de $\lambda$.\\
\textbf{El truncamiento no toca la magnitud que se estima.}}
\end{minipage}};
\end{tikzpicture}
\end{center}

\vspace{2mm}
\centering{\footnotesize\color{gris}\textbf{Salvedad declarada:} el
adelgazamiento no es independiente de la posición si las instalaciones grandes
se agrupan de forma distinta a las pequeñas. La correlación de marcas es el
diagnóstico de eso, y no muestra tal patrón.}

\end{frame}

% #############################################################################
\section{La escala de influencia}
% #############################################################################

\begin{frame}{Superficie de presión de emisión}
\vspace{-10mm}
\begin{columns}[T]
\begin{column}{0.55\textwidth}
\fig[1.5]{28_presion_emision.png}
\end{column}
\begin{column}{0.4\textwidth}
\[
P_h(\loc) = \sum_{i} \log(1+m_i)\,
            \kappa_h\!\left(\lVert\loc-\loci{i}\rVert\right)
\]

\vspace{4mm}
{\small\textbf{¿Por qué $\log(1+m_i)$ y no $m_i$?}}

\vspace{2mm}
{\small La asimetría de las marcas es \num{3.30}: las diez mayores
instalaciones concentran el \SI{17.2}{\percent} del \ce{NOx} total.}

\vspace{2mm}
{\small Sin transformar, unas pocas centrales térmicas dominarían por completo
la superficie.}
\end{column}
\end{columns}

\end{frame}

% =============================================================================
\begin{frame}{La escala se estima, no se fija por diseño}
\vspace{-5mm}
\begin{columns}[T]
\begin{column}{0.5\textwidth}
\fig[1.1]{30_escala_influencia.png}
\end{column}
\begin{column}{0.46\textwidth}
\vspace{1mm}
\centering
\small
\begin{tabular}{@{}lcc@{}}
\toprule
$h$ (\si{\kilo\metre}) & RMSE & $\Delta$AIC \\
\midrule
--- (base) & \num{3.657} & \num{2.777} \\
\num{5} & \num{3.688} & \num{2.777} \\
\rowcolor{grisc}
\textbf{10} & \textbf{\num{3.642}} & \textbf{\num{0.000}} \\
\num{15} & \num{3.654} & \num{1.463} \\
\num{25} & \num{3.669} & \num{5.800} \\
\num{40} & \num{3.666} & \num{7.121} \\
\num{100} & \num{3.674} & \num{12.937} \\
\bottomrule
\end{tabular}

\vspace{3mm}
{\small Curva en U limpia: peor a \num{5} y a \num{15}.}

\vspace{3mm}
{\footnotesize $F = \num{13.86}$, \; $p = \num{2.5e-4}$\\
$R^2$ ajustado \num{0.754} $\rightarrow$ \num{0.769}}
\end{column}
\end{columns}

\end{frame}

% =============================================================================
\begin{frame}{Lo que este resultado \emph{no} dice}
\vspace{-10mm}
\vspace{4mm}
\begin{columns}[T]
\begin{column}{0.46\textwidth}
\begin{block}{Lo que sí}
\small
\begin{itemize}
  \item Efecto estadísticamente robusto ($p = \num{2.5e-4}$)
  \item Escala identificada con curva en U limpia
  \item Multiplicar por $e$ la presión eleva el \ce{NO2} un
        \SI{6.35}{\percent}
\end{itemize}
\end{block}
\end{column}
\begin{column}{0.46\textwidth}
\begin{alertblock}{Lo que no}
\small
\begin{itemize}
  \item La mejora del RMSE es del \textbf{\SI{0.41}{\percent}}
  \item Poder predictivo escaso
  \item Coherente con Canarias fuera del dominio (\SI{24}{\percent} del
        \ce{NOx} industrial)
\end{itemize}
\end{alertblock}
\end{column}
\end{columns}

\vspace{7mm}
\begin{center}
{\small Declararlo aquí y no enterrarlo en las limitaciones: significación
estadística y relevancia predictiva son cosas distintas, y confundirlas es el
error más común al leer un $p$-valor.}
\end{center}

\end{frame}

% =============================================================================
\begin{frame}{El hallazgo central}

\vspace{2mm}
\begin{center}
\small
\begin{tabular}{@{}p{4.3cm}p{4.4cm}c@{}}
\toprule
\textbf{Magnitud} & \textbf{Datos usados} & \textbf{km} \\
\midrule
Diámetro del distrito\newline
{\footnotesize (Thomas, rama B)} &
{\footnotesize solo coordenadas de instalaciones} & \num{10.7} \\[3mm]
Escala de influencia\newline
{\footnotesize (barrido, rama C)} &
{\footnotesize coordenadas + concentraciones} & \num{10.0} \\
\midrule
\rowcolor{grisc}
\multicolumn{2}{@{}l}{\textbf{Razón}} & {\color{rojo}\textbf{\num{0.93}}} \\
\bottomrule
\end{tabular}
\end{center}

\vspace{3mm}
\begin{block}{}
\centering\large
La unidad de influencia sobre el \ce{NO2} es el \textbf{distrito industrial},
no la instalación aislada.
\end{block}

\vspace{1mm}
{\scriptsize\color{gris}\textbf{Robustez:} la convergencia se
verificó dos veces sobre patrones distintos. Antes de depurar el preprocesado
—que colapsó tres instalaciones duplicadas— \num{10.0}/\num{9.5} =
\num{1.05}; después, \num{0.93}. Sobrevivir a un cambio no planeado en los
datos de entrada es mejor evidencia que cualquier remuestreo.}

\end{frame}

% =============================================================================
% =============================================================================
% LA MISMA ESCALA, SOBRE EL MAPA. Va después de "El hallazgo central" porque
% el tercer panel usa el diámetro de 10,7 km que allí se establece.
\begin{frame}{Esa escala, frente a las unidades en que se legisla}

\vspace{-2mm}
\begin{center}
\includegraphics[width=1\textwidth]{45_tres_escalas.png}
\end{center}

\vspace{-6mm}
\begin{columns}[T]
\begin{column}{0.31\textwidth}
\centering
\textbf{\SI{31158}{\kilo\metre\squared}}\\
{\footnotesize por comunidad}
\end{column}
\begin{column}{0.31\textwidth}
\centering
\textbf{\SI{9970}{\kilo\metre\squared}}\\
{\footnotesize por provincia}
\end{column}
\begin{column}{0.34\textwidth}
\centering
{\color{rojo}\textbf{\SI{3.6}{\percent} del territorio}}\\
{\footnotesize \num{169} distritos, $\varnothing$ \SI{10.7}{\kilo\metre}}
\end{column}
\end{columns}

\vspace{-2mm}
\begin{center}
{\small El mecanismo vive en el \SI{3.6}{\percent} del territorio; las unidades
sobre las que se promedia son dos órdenes de magnitud mayores.}
\end{center}

\vspace{-3mm}
\centering{\scriptsize\color{gris}Los distritos son la unión de discos del
diámetro estimado en torno a las \num{297} instalaciones: la huella del
mecanismo, no los \num{520} centros latentes.}

\end{frame}

% =============================================================================
\begin{frame}{No es falta de potencia. Es la unidad de análisis.}

\begin{columns}[T]
\begin{column}{0.44\textwidth}
\vspace{-2mm}
\begin{center}
\begin{tikzpicture}[scale=1]
  \draw[fill=gray!12, draw=gray!50] (0,0) circle (2.5);
  \node[gris] at (0,1.95) {\footnotesize provincia media};
  \node[gris] at (0,1.6) {\footnotesize \SI{9970}{\kilo\metre\squared}};
  \node[gris] at (0,-2.05) {\footnotesize $\varnothing$ equiv.\
        \SI{113}{\kilo\metre}};
  \fill[rojo] (0,0) circle (0.22);
  \node[rojo,right] at (0.42,0.22) {\footnotesize diámetro \SI{10}{\kilo\metre}};
  \node[rojo,right] at (0.42,-0.18) {\footnotesize\textbf{\SI{0.8}{\percent}
        del área}};

\end{tikzpicture}
\end{center}
\end{column}
\begin{column}{0.52\textwidth}
\vspace{4mm}
\fig[1.1]{50_curva_escala.png}
\end{column}
\end{columns}

\vspace{2mm}
\begin{center}
\large
La urbanización \textbf{opera} a escala provincial. La industria no.
\end{center}
\vspace{1mm}
\begin{center}
{\small La diferencia no está en la intensidad del mecanismo sino en su
\textbf{alcance espacial}.}
\end{center}

\end{frame}

% #############################################################################
\section{Conclusiones}
% #############################################################################

\begin{frame}{Conclusiones}

\vspace{2mm}
\begin{enumerate}
  \item[\textbf{1}] Dos ramas metodológicamente independientes sitúan la
  escala industrial en \textbf{$\approx$\SI{10}{\kilo\metre}} (razón
  \num{0.93}). La unidad de influencia es el \textbf{distrito}, no la
  instalación aislada.

  \vspace{3mm}
  \item[\textbf{2}] El mismo efecto es \textbf{indetectable} a escala
  provincial. No por falta de potencia: porque la unidad de análisis promedia
  la señal sobre un territorio dos órdenes de magnitud mayor que su alcance.

  \vspace{3mm}
  \item[\textbf{3}] El MAUP pasa de advertencia genérica a
  \textbf{diagnóstico cuantitativo} cuando se dispone de una escala de
  referencia estimada externamente.
\end{enumerate}

\vspace{3mm}
\begin{center}
\begin{tikzpicture}
\node[draw=azul,line width=0.7pt,rounded corners=2pt,inner sep=3.5mm,
      fill=grisc]{%
\begin{minipage}{0.78\textwidth}\centering
{\color{azul}Sin la rama de procesos puntuales, el resultado nulo areal sería
indistinguible de una ausencia real de efecto.}
\end{minipage}};
\end{tikzpicture}
\end{center}

\end{frame}

% =============================================================================
\begin{frame}{Limitaciones y trabajo futuro}

\begin{columns}[T]
\begin{column}{0.47\textwidth}
\textbf{Limitaciones}
\begin{itemize}\footnotesize
  \item \textbf{Umbral E-PRTR de \SI{100}{\tonne}}: solo grandes emisores;
  sesga $\lambda$, no $K$.
  \item \textbf{Canarias fuera del dominio}, con el \SI{24}{\percent} del
  \ce{NOx} industrial español.
  \item \textbf{Se modela fondo, no exposición máxima}: los picos de tráfico
  ocurren por debajo de la resolución de \SI{2}{\kilo\metre}.
  \item \textbf{CORINE clase 121} mezcla industria con comercio y logística.
  \item \textbf{$n = \num{50}$ unidades areales}: poca potencia.
  \item \textbf{Envolventes con $\lambda$ fijada} de los propios datos:
  contraste anticonservador.
\end{itemize}
\end{column}

\begin{column}{0.47\textwidth}
\textbf{Trabajo futuro}
\begin{itemize}\footnotesize
  \item Formulación bayesiana con \textbf{\textsc{spde}} para propagar la
  incertidumbre de la superficie predicha hacia el modelo areal, en lugar de
  tratarla como dato.
  \item Extensión temporal \textbf{2019--2023} explotando cierres de
  instalaciones como cuasiexperimento.
  \item Repetir el variograma con modelo esférico como análisis de
  sensibilidad.
\end{itemize}

\vspace{3mm}
\begin{block}{Reproducibilidad}
\footnotesize
\texttt{R/99\_verificar.R} comprueba invariantes de conservación y frescura de
los ficheros derivados. Estado actual: \textbf{0 fallos}.
\end{block}
\end{column}
\end{columns}

\end{frame}

% =============================================================================
\begin{frame}[plain]
\vfill
\hspace*{8mm}\begin{minipage}{0.85\textwidth}
  {\color{azulc}\rule{\textwidth}{1.2pt}}\\[6mm]
  {\Large\color{azul}\bfseries Gracias}\\[7mm]
  {\small Todo el trabajo es reproducible y está publicado:}\\[2mm]
  {\normalsize\color{azul}\texttt{github.com/randyfernandezp/no2-espana}}\\[2mm]
  {\footnotesize\color{gris}Código de descarga, análisis y figuras, más
  \texttt{R/99\_verificar.R}, que comprueba la integridad de la cadena.}\\[7mm]
  {\color{azulc}\rule{\textwidth}{0.5pt}}\\[2mm]
  {\footnotesize\color{gris}Randy Fernández Palomino \textperiodcentered\
  José Carhuaricra Cusipuma \hfill 2026}
\end{minipage}
\vfill
\end{frame}

% #############################################################################
% DIAPOSITIVAS DE RESERVA — para preguntas del tribunal
% #############################################################################
\appendix

\begin{frame}{Reserva: siglas}
\vspace{-10mm}\begin{columns}[T]
\begin{column}{0.49\textwidth}
\scriptsize
\begin{tabular}{@{}ll@{}}
\toprule
\multicolumn{2}{@{}l}{\textbf{Método}} \\
\midrule
\textsc{ace} & aleatoriedad espacial completa \\
 & (\emph{complete spatial randomness}) \\
\textsc{maup} & problema de la unidad de área modificable \\
 & (\emph{modifiable areal unit problem}) \\
\textsc{lisa} & indicadores locales de asociación espacial \\
\textsc{mco} & mínimos cuadrados ordinarios \\
\textsc{sar} & autorregresivo simultáneo \\
\textsc{car} & autorregresivo condicional \\
\textsc{sem} & modelo de error espacial \\
\textsc{gmrf} & campo markoviano aleatorio gaussiano \\
\textsc{spde} & ecuación diferencial estocástica en \\
 & derivadas parciales \\
\textsc{ko}, \textsc{kde} & kriging ordinario y con deriva externa \\
\bottomrule
\end{tabular}
\end{column}
\begin{column}{0.49\textwidth}
\scriptsize
\begin{tabular}{@{}ll@{}}
\toprule
\multicolumn{2}{@{}l}{\textbf{Diagnóstico y datos}} \\
\midrule
\textsc{rmse} & raíz del error cuadrático medio \\
\textsc{aic} & criterio de información de Akaike \\
\textsc{vif} & factor de inflación de la varianza \\
\textsc{lmg} & descomposición de Lindeman, Merenda \\
 & y Gold del $R^2$ \\
\midrule
\textsc{nuts} & nomenclatura de unidades territoriales \\
 & estadísticas (Eurostat) \\
\textsc{eea} & Agencia Europea de Medio Ambiente \\
\textsc{ied} & Directiva de Emisiones Industriales \\
\textsc{e-prtr} & registro europeo de emisiones \\
 & y transferencias de contaminantes \\
\textsc{corine} & inventario europeo de cobertura del suelo \\
\textsc{lur} & regresión de uso del suelo \\
\bottomrule
\end{tabular}
\end{column}
\end{columns}
\end{frame}

\begin{frame}{Reserva: notación}
\vspace{-5mm}
\begin{columns}[T]
\begin{column}{0.49\textwidth}
\centering
\scriptsize
\begin{tabular}{@{}ll@{}}
\toprule
\multicolumn{2}{@{}l}{\textbf{Ramas continua y areal}} \\
\midrule
$\dom$ & área de análisis \\
$\loci{i}$ & localización $i$-ésima (no aleatoria) \\
$\dvec$, \; $\dsc$ & desplazamiento y distancia \\
$\Wp{\loc}$, \; $\Yp{\loc}$ & proceso latente y observable \\
$\varesp$, \; $\pepita$ & varianza espacial y pepita \\
$\escala$, \; $\suave$ & escala y suavizado de la correlación \\
$\rangop$ & rango práctico: $\corr(\rangop)=\num{0.05}$ \\
\midrule
$\{R_1,\dots,R_{50}\}$ & índice discreto (provincias) \\
$\Adj$, \; $\Prec$ & adyacencia y precisión \\
$\alpha$ & parámetro autorregresivo, $|\alpha|<1$ \\
\bottomrule
\end{tabular}
\end{column}
\begin{column}{0.49\textwidth}
\centering
\scriptsize
\begin{tabular}{@{}ll@{}}
\toprule
\multicolumn{2}{@{}l}{\textbf{Rama de procesos puntuales}} \\
\midrule
$N(A)$ & nº de eventos en $A$ \\
$\lambda(\loc)$ & intensidad de primer orden \\
$\lambda_2(\loc,\loc')$ & intensidad de segundo orden \\
$g_2$, \; $K(t)$ & correlación por pares y función $K$ \\
$\Lambda(\loc)$ & intensidad aleatoria (Cox) \\
$\kappa$, \; $\sigma$ & padres y dispersión \\
$h$, \; $p$ & ancho de banda y prob.\ de retención \\
$Z_j(\loc)$ & covariable del modelo de intensidad \\
\bottomrule
\end{tabular}
\end{column}
\end{columns}

\vspace{5mm}
{\footnotesize\color{gris} Las localizaciones son $\loc$, nunca $x/y$; el
desplazamiento, $\dvec$. La adyacencia es $\Adj$: la letra $W$ está reservada
al proceso. $\rho$ designa \emph{solo} la correlación espacial: el
autorregresivo es $\alpha$ y la intensidad de padres, $\kappa$.}
\end{frame}

\begin{frame}{Reserva: cómo se define la vecindad}
\begin{columns}[T]
\begin{column}{0.54\textwidth}
\fig[0.8]{44_grafo_vecindad.png}
\end{column}
\begin{column}{0.42\textwidth}
\vspace{4mm}
{\small El grafo no dirigido de contigüidad tiene \textbf{cuatro componentes
conexas}: península, Mallorca, Menorca y Eivissa--Formentera.}

\vspace{3mm}
{\small Las islas quedan sin vecinos con quien compartir información y los
contrastes de dependencia no están definidos sobre un grafo disconexo.}

\vspace{4mm}
\begin{block}{}
\small
Con $k=5$ vecinos más próximos el grafo se reconecta. Es la razón de usar
$k$-vecinos y no contigüidad en todo el diagnóstico areal.
\end{block}
\end{column}
\end{columns}
\end{frame}

\begin{frame}{Reserva: ¿hay anisotropía?}
\vspace{-10mm}
\begin{columns}[T]
\begin{column}{0.52\textwidth}
\fig[1.2]{03_variograma_direccional.png}
\end{column}
\begin{column}{0.44\textwidth}
\vspace{6mm}
{\small Variogramas en cuatro direcciones. El semivariograma máximo varía
entre \num{0.87} y \num{2.09} según la dirección.}

\vspace{4mm}
{\small Diferencias de esa magnitud sugieren anisotropía zonal, no corregida
en este trabajo. Es una limitación adicional del modelo de campo, y una vía
directa de mejora.}
\end{column}
\end{columns}
\end{frame}

\begin{frame}{Reserva: ¿las instalaciones próximas emiten cantidades parecidas?}
\vspace{-10mm}
\begin{columns}[T]
\begin{column}{0.52\textwidth}
\fig[1.5]{27_markcorr.png}
\end{column}
\begin{column}{0.44\textwidth}
\vspace{6mm}
{\small Función de correlación de marcas.}

\vspace{4mm}
{\small Si fuera plana, la masa emitida se distribuiría con independencia de
la posición, y bastaría el patrón sin marcar para construir la superficie de
presión.}

\vspace{4mm}
{\small Es la justificación empírica de ponderar por $\log(1+m_i)$ en la
ecuación de presión.}
\end{column}
\end{columns}
\end{frame}

\begin{frame}{Reserva: probabilidad de superar umbrales}
\begin{columns}[T]
\begin{column}{0.48\textwidth}
\centering
{\footnotesize $P(\ce{NO2} > \SI{20}{\ugm})$ --- directiva 2030}\\[1mm]
\fig[1.2]{10_prob_20.png}
\end{column}
\begin{column}{0.48\textwidth}
\centering
{\footnotesize $P(\ce{NO2} > \SI{10}{\ugm})$ --- recomendación OMS}\\[1mm]
\fig[1.2]{12_prob_10.png}
\end{column}
\end{columns}
\vspace{-1mm}
\centering{\small Superficie nacional con $P > 0{,}5$: \SI{0.16}{\percent}
para el umbral de \SI{20}{\ugm}. Población con $P > 0{,}5$ de superar
\SI{10}{\ugm}: \SI{9.26}{\percent}.}
\end{frame}

\begin{frame}{Reserva: residuos del modelo de intensidad}
\vspace{-5mm}
\begin{columns}[T]
\begin{column}{0.52\textwidth}
\fig[1.5]{25_residuos_ppm.png}
\end{column}
\begin{column}{0.44\textwidth}
\vspace{6mm}
{\small Residuos suavizados del proceso de Poisson inhomogéneo.}

\vspace{4mm}
{\small Zonas positivas indican más instalaciones de las que predicen las
covariables; negativas, menos.}

\vspace{4mm}
{\small Es el diagnóstico que motiva pasar a un modelo de conglomerados.}
\end{column}
\end{columns}
\end{frame}

\end{document}
